\section{Related Work}
\label{sec:related}

\subsection{Sequential Recommendation}
SASRec~\cite{kang2018self} and BERT4Rec~\cite{sun2019bert4rec} set the
benchmark for next-item prediction using causal and bidirectional
self-attention over interaction histories.
Both are ID-based: each item is a learned lookup embedding, so unseen
items cause cold-start failure without retraining.
DIF-SASRec~\cite{xie2022difsr} partially mitigates this with decoupled
attention streams for IDs, attributes, and content, but auxiliary features
remain supplementary to the ID embedding.
All three models train offline, do not generalize to unseen IDs, and lack
support for multilingual queries or real-time adaptation.

\subsection{Dense and Multimodal Retrieval}
CLIP~\cite{radford2021clip} aligns images and text in a shared latent
space via large-scale contrastive pretraining, but queries are encoded
independently with no user history, and its English corpus limits
multilingual use.
BGE-M3~\cite{chen2024bgem3} extends dense retrieval to multilingual,
multi-granularity queries with both symmetric and asymmetric matching,
but it is similarly stateless and provides no behavioral signal.

\subsection{Graph-Based Collaborative Filtering}
NGCF~\cite{wang2019ngcf}, LightGCN~\cite{he2020lightgcn}, and
PinSage~\cite{ying2018pinsage} learn from co-purchase edges via graph
convolutions or random-walk sampling.
GNNs capture high-order connectivity but aggregate over the full
interaction graph, so memory and compute costs grow with catalog size
and GPU hardware is typically required.
Cleora~\cite{cleora2021} avoids factorization and GPU training through
deterministic Markov-chain propagation, but still needs at least one
co-purchase edge per item. Graph methods are also query-agnostic: they return
the same neighborhood regardless of the typed query.

\subsection{Hybrid Retrieval Systems and Remaining Gaps}
Existing approaches — collaborative filtering~\cite{wang2019ngcf,he2020lightgcn,ying2018pinsage},
content-based retrieval~\cite{radford2021clip}, and sequential
models~\cite{kang2018self,xie2022difsr} — share three structural limitations.
First, they use a single pipeline with no separation between
user-initiated search and system-initiated suggestions.
Second, most fuse at most two modalities and assume a static user
representation not updated between requests.
Third, user state is maintained implicitly through training rather than
as an explicitly persisted embedding updated after each interaction.
ID-based sequential models cannot generalize to unseen items without
retraining. Dense retrievers return relevant results but lack user
identity or interaction history. Graph methods capture co-purchase
signals but cover a fraction of the catalog and ignore query content.
No existing system simultaneously supports multilingual and image-based
active search, graph-based and sequential passive recommendations, and
a continuously-updated shared user profile.

\subsection{Text Encoders for Item Retrieval}
BLaIR~\cite{hou2024amazon} fine-tunes a RoBERTa-Large backbone on
(review text, item) contrastive pairs, producing embeddings aligned
with review-style queries.
Two limitations drive our encoder choice.
First, BLaIR is English-only and degrades on short keyword queries
outside the review-text distribution.
Second, an evaluation-bias problem arises when benchmark clicks are
collected from a system using the same encoder: held-out items are
correlated with that encoder, inflating sampled-evaluation metrics.
We address this in Section~\ref{ssec:encoder} by supplementing
sampled evaluation with live retrieval quality (genre-precision@10
over 22 labeled queries) where no circular bias applies.
